% !TEX root = task-2-11-1.tex
% Editable task bank. Append new BankTask blocks at the end of the document.
\documentclass[12pt,a4paper]{article}
\usepackage{physics-textbook}
\usepackage{graphicx,needspace}
\usetikzlibrary{decorations.pathmorphing}
\newif\ifShowSolutions
\ifdefined\TutorEdition
  \ShowSolutionstrue
\else
  \ShowSolutionsfalse % Change to \ShowSolutionstrue to display solutions.
\fi
\newcounter{banktask}
\newsavebox{\BankTaskBox}
\newcommand{\BankTaskContent}[3]{%
  \noindent\textbf{\thebanktask.}\quad #2\par
  \ifShowSolutions
    \Needspace{50mm}% Keep the solution heading with its opening text or diagram.
    #3\par
    \ifdefined\TutorEdition
      \ifcsname TutorSource#1\endcsname
        \nopagebreak\smallskip{\small\csname TutorSource#1\endcsname\par}%
      \fi
    \fi
  \fi
}
\newcommand{\BankTask}[3]{%
  \refstepcounter{banktask}\par\medskip
  \sbox{\BankTaskBox}{\begin{minipage}{\linewidth}%
    \BankTaskContent{#1}{#2}{#3}\end{minipage}}%
  \ifdim\dimexpr\ht\BankTaskBox+\dp\BankTaskBox\relax>\dimexpr\textheight-20mm\relax
    \Needspace{35mm}\BankTaskContent{#1}{#2}{#3}%
  \else
    \Needspace{\dimexpr\ht\BankTaskBox+\dp\BankTaskBox\relax}%
    \noindent\usebox{\BankTaskBox}\par
  \fi
  \medskip
}
\begin{document}
\BookSubsection{2.11.1}{Работа газа в изобарном процессе}


% Source sdamgia-37890: https://phys-ege.sdamgia.ru/problem?id=37890
\BankTask{sdamgia-37890}{
При сжатии идеального одноатомного газа при постоянном давлении внешние силы совершили работу $800\,\text{Дж}$. Чему равно изменение внутренней энергии газа? (В джоулях.)
}{
\smallskip\noindent\textbf{Решение}\par
Работа газа при изобарном процессе равна $A=\nu R\Delta T= - 800\text{Дж},$ изменение внутренней энергии одноатомного идеального газа при изобарном процессе: 

$\Delta U = \frac{3}{2} \nu R\Delta T = \frac{3}{2} A = \frac{3}{2} \cdot ( - 800 ) = - 1200 \text{Дж}.$

Ответ: $-1200$.
}

% Source sdamgia-12858: https://phys-ege.sdamgia.ru/problem?id=12858
\BankTask{sdamgia-12858}{
В сосуде объемом $250\,\text{л}$ под тяжелым поршнем находится кислород. Давление кислорода $300\,\text{кПа}$. В изобарном процессе плотность газа уменьшилась в $5$ раз. Какую работу совершил газ в этом процессе? (В килоджоулях.)
}{
\smallskip\noindent\textbf{Решение}\par
Поскольку плотность газа уменьшилась в $5$ раз, его объем увеличился в $5$ раз и стал $5\cdot 250=1250\,\text{л}$. Работа газа в изобарном процессе равна

$A=p\Delta V=300 \cdot 10 ^{3} \cdot ( 1250 - 250 ) \cdot 10 ^{- 3} =300 \cdot 10 ^{3} \text{Дж}=300\text{кДж}.$

Ответ: $300$.
}

% Source sdamgia-3517: https://phys-ege.sdamgia.ru/problem?id=3517
\BankTask{sdamgia-3517}{
Для анализа изотермического, изобарного и изохорного процессов над фиксированным количеством идеального газа используют первое начало термодинамики: $Q=A + \Delta U.$ Передаваемое количество теплоты при:

\TaskMatchingColumns{Изопроцесс}{А)  Изотермическом процессе\par Б)  Изобарном процессе\par В)  Изохорном процессе}{Физическое явление}{$1$.  Идет на увеличение его внутренней энергии\par $2$.  Полностью превращается в работу\par $3$.  Идет на увеличение его внутренней энергии и на работу}

\begin{center}\begin{tabular}{|c|c|c|}\hline
$A$ & Б & В\\\hline
\strut & \strut & \strut\\\hline\end{tabular}\end{center}
}{
\smallskip\noindent\textbf{Решение}\par
Вспомним все, что может потребоваться при решении данной задачи. 

Для идеального газа внутренняя энергия зависит только от температуры: $\Delta U= \frac{3}{2} \nu R\Delta T.$ Следовательно, изменить внутреннюю энергию идеального газа можно, только нагрев или охладив его. Работа газа всегда связана с изменением его объема, если газ расширяется, он совершает положительную работу, если сжимается, то отрицательную. Наконец, идеальный газ подчиняется уравнению состояния Клапейрона — Менделеева: $pV=\nu RT.$

При изотермическом процессе температура остается неизменной, поэтому внутренняя энергия идеального газа не изменяется, а значит, согласно первому началу термодинамики, все переданное тепло полностью превращается в работу (А — $2$). При изобарном процессе, то есть при постоянном давлении, согласно уравнению Клапейрона — Менделеева, меняется как температура, так и объем газа, а значит, в правой части первого начала термодинамики отличны от нуля оба слагаемых: тепло идет и на увеличение внутренней энергии газа, и на работу (Б — $3$). Наконец, при изохорном процессе объем газа фиксирована. Следовательно, все тепло идет на изменение внутренней энергии (В — $1$).

Ответ: $231$.
}

% Source sdamgia-3107: https://phys-ege.sdamgia.ru/problem?id=3107
\BankTask{sdamgia-3107}{
Используя первый закон термодинамики, установите соответствие между особенностями теплового процесса в идеальном газе и его названием. 

\TaskMatchingColumns{Особенности теплового процесса}{А)  Все передаваемое газу количество теплоты идет на совершение газом работы\par Б)  Изменение внутренней энергии газа равно количеству переданной теплоты, при этом газ не совершает работы}{Название теплового процесса}{$1$.  Изохорный\par $2$.  Изотермический\par $3$.  Изобарный\par $4$.  Адиабатный}

\begin{center}\begin{tabular}{|c|c|}\hline
$A$ & Б\\\hline
\strut & \strut\\\hline\end{tabular}\end{center}
}{
\smallskip\noindent\textbf{Решение}\par
Согласно первому началу термодинамики, переданное газу тепло идет на изменение его внутренней энергии и на совершение газом работы против внешних сил: $Q=\Delta U + A.$ В процессе $A$ все передаваемое идеальному газу количество теплоты идет на совершение газом работы, а значит, внутренняя энергия газа не изменяется. Следовательно, не изменяется температура газа, процесс $A$ — изотермический (А — $2$). В процессе $\text{Б}$ газ не совершает работы. Это означает, что его объем не изменяется. Такой процесс называется изохорным (Б — $1$).

Ответ: $21$.
}

% Source sdamgia-9326: https://phys-ege.sdamgia.ru/problem?id=9326
\BankTask{sdamgia-9326}{
При сжатии идеального одноатомного газа при постоянном давлении внешние силы совершили работу $2000\,\text{Дж}$. Какое количество теплоты было передано при этом газом окружающим телам? (Ответ дайте в джоулях.)
}{
\smallskip\noindent\textbf{Решение}\par
Изменение внутренней энергии идеального одноатомного газа при постоянном давлении равно $\Delta U= \frac{3}{2} \nu R\Delta T,$ работа газа при изобарном процессе $A=\nu R\Delta T.$ Причем, изменение внутренней энергии при изобарном сжатии отрицательное:

\[\Delta U= - \frac{3}{2} A_{\,\text{внеш}}.\]

Количество теплоты, отданное газом окружающим телам, по первому закону термодинамики равно:

\[Q_{\,\text{отд}}=A_{\,\text{внеш}} - \Delta U=A_{\,\text{внеш}} + \frac{3}{2} A_{\,\text{внеш}}= \frac{5}{2} A_{\,\text{внеш}}= \frac{5}{2} \cdot 2000 = 5000\,\text{Дж}.\]

Ответ: $5000\,\text{Дж}$.
}

% Source sdamgia-9753: https://phys-ege.sdamgia.ru/problem?id=9753
\BankTask{sdamgia-9753}{
Неизменное количество идеального одноатомного газа изохорически переводят из состояния $1$ в состояние $2$. Затем газ изобарически переводят из состояния $2$ в состояние $3$, и при этом газ совершает работу $40\,\text{Дж}$. Известно, что температура газа в процессе $2$-$3$ повышается на столько же, на сколько она повысилась в процессе $1$-$2$. Какое количество теплоты поглотил газ в процессе $1$-$2$?
}{
\smallskip\noindent\textbf{Решение}\par
В изобарном процессе совершается работа $A_{23}=\nu R\Delta T.$ 

В изохорном процессе работа газа равна нулю, и все тепло идет исключительно на увеличение внутренней энергии:

$Q_{12}=\Delta U_{12}= \frac{3}{2} \nu R\Delta T= \frac{3}{2} A_{23}= \frac{3}{2} \cdot 40 =60\text{Дж}.$

Ответ: $60\,\text{Дж}$.
}

% Source sdamgia-11893: https://phys-ege.sdamgia.ru/problem?id=11893
\BankTask{sdamgia-11893}{
Постоянную массу идеального одноатомного газа изобарно сжали так, что $T_2 = \frac{T_1}{k} .$ Затем этот же газ адиабатически расширяется так, что $T_3 = \frac{T_2}{2} .$ Отношение модулей работ в изобарном и адиабатическом процессах $n = 4.$ Найдите $k$.
}{
\smallskip\noindent\textbf{Решение}\par
Работа, совершенная над газом при изобарном сжатии, вычисляется по формуле: $A_1 = p|\Delta V|.$ Из уравнения Менделеева — Клапейрона для первого и второго состояния газа $p_1V_1 = \nu RT_1$ и $p_2V_2 = \nu RT_2$ получаем:

$A_1 = p ( V_1 - V_2 ) = \nu R ( T_1 - T_2 ) = \nu R ( T_1 - \frac{T_1}{k} ) = \frac{k - 1}{k} \nu RT_1.$

Согласно первому началу термодинамики при адиабатическом расширении работа газа равна $A_2 = - \Delta U,$ откуда:

$A_2 = - \frac{3}{2} \nu R ( T_3 - T_2 ) = \frac{3}{2} \nu R ( T_2 - T_3 ) = \frac{3}{4} \nu R T_2 = \frac{3}{4k} \nu RT_1.$

Отношение работ равно $4$:

$\frac{A_1}{A_2} = 4,$ откуда $\frac{k - 1}{k} \cdot \frac{4k}{3} =4,$ следовательно, $k=4$.

Ответ: $4$.
}

% Source sdamgia-11953: https://phys-ege.sdamgia.ru/problem?id=11953
\BankTask{sdamgia-11953}{
Гелий в количестве $\nu =3\,\text{моль}$ изобарно сжимают, совершая работу $A_{1}=2{,}4\,\text{кДж}$. При этом температура гелия уменьшается в $4$ раза: $T_2 = \frac{T_1}{4} .$ Затем газ адиабатически расширяется, при этом его температура изменяется до значения $T_3 = \frac{T_1}{8} .$ Найдите работу газа $A_{2}$ при адиабатном расширении. Количество вещества в процессах остается неизменным.
}{
\smallskip\noindent\textbf{Решение}\par
При изобарном сжатии над гелием совершается работа, модуль которой $A_1 = |p\Delta V|,$ где $\text{р}$ — давление гелия в этом процессе, $\Delta V$ — изменение его объема.

В соответствии с уравнением Менделеева—Клапейрона для этого процесса можно записать:

$p\Delta V = \nu R ( T_1 - T_2 ) .$

В адиабатном процессе (процессе без теплообмена) в соответствии с первым законом термодинамики сумма изменения внутренней энергии газа $\Delta U_2= \frac{3}{2} \nu R ( T_3 - T_2 )$ и его работы $A_2$ равна $\Delta U_2 + A_2=0.$ Откуда получаем:

$\frac{3}{2} \nu R ( T_3 - T_2 ) + A_2 = 0.$

Преобразуя записанные уравнения с учетом соотношений температур, заданных в условии задачи, получаем:

$A_1 = 3\nu RT_2;$ $A_2 = \frac{3}{4} \nu R T_2.$

Тогда:

$A_2 = \frac{A_1}{4} = \frac{2400\text{Дж}}{4} = 600\text{Дж}.$

Ответ: $600\,\text{Дж}$.
}

% Source sdamgia-25410: https://phys-ege.sdamgia.ru/problem?id=25410
\BankTask{sdamgia-25410}{
В вертикальном цилиндре, закрытом легким поршнем, находится бензол ($\mathrm{C_6H_6}$) при температуре кипения t  =  $80$ \ensuremath{{}^\circ}C. При сообщении бензолу количества теплоты Q часть его превращается в пар, который при изобарном расширении совершает работу А. Удельная теплота парообразования бензола L  =  $396$ · $10^{3}$ Дж/кг, его молярная масса M  =  $78$ · $10^{-3}$ кг/моль. Какая часть подведенного к бензолу количества теплоты переходит в работу? Объемом жидкого бензола пренебречь.
}{
\smallskip\noindent\textbf{Решение}\par
$1$.   При кипении бензола происходит его изобарное расширение. Работа пара $A = p\Delta V,$ где p  — атмосферное давление, $\Delta V$  — изменение объема.
$2$.  Считая пар идеальным газом, воспользуемся уравнением Клапейрона  — Менделеева для определения изменения объема за счет испарившегося бензола массой $\Delta m:$ $p\Delta V= \frac{\Delta m}{M} RT,$ где M  =  $78$ · $10^{-3}$ кг/моль  — молярная масса бензола, T  =  $80$ + $273$  =  $353$ K  — температура кипения бензола. Отсюда $A= \frac{\Delta mRT}{M} .$
$3$.  Количество теплоты Q, необходимое для испарения массы $\Delta m$ бензола, пропорционально удельной теплоте парообразования L: $Q = \Delta mL .$
$4$.  Искомая величина определяется отношением:

\[\eta = \frac{A}{Q} = \frac{RT}{ML} = \frac{8{,}31 \cdot 353}{78 \cdot 10 ^{ - 3 } \cdot 396 \cdot 10 ^3} \approx 0{,}095.\]

Ответ:$\eta \approx 0{,}095.$
}

% Source sdamgia-33046: https://phys-ege.sdamgia.ru/problem?id=33046
\BankTask{sdamgia-33046}{
В вертикальном цилиндре с гладкими стенками под массивным металлическим поршнем находится одноатомный идеальный газ. В начальном состоянии поршень массой M и площадью основания S покоится на высоте h, опираясь на выступы (см. рис. $1$). Давление газа $p_{0}$ равно внешнему атмосферному. Какое количество теплоты Q нужно сообщить газу при медленном его нагревании, чтобы поршень оказался на высоте H (cм. рис. $2$)? Тепловыми потерями пренебречь.

\begin{center}\begin{tikzpicture}[task graph,>=stealth]\draw (0,3.5)--(0,0)--(2.25,0)--(2.25,3.5);\draw(-.1,0)--(2.4,0);\foreach \x in {0,.16,...,2.4}{\draw(\x,0)--++(-.15,-.16);}\fill(0,1.2)--(.23,1.2)--(0,.97)--cycle;\fill(2.25,1.2)--(2.02,1.2)--(2.25,.97)--cycle;\draw[fill=orange!30](0,1.2)rectangle++(2.25,.32);\draw(2.25,1.2)--(3.1,1.2);\draw[<->](2.85,0)--(2.85,1.2)node[midway,right]{\(h\)};\node at(1.05,2.45){\(p_0\)};\node at(1.1,.53){\(p_0\)};\draw(1.8,1.52)--(2.65,2.2)node[right]{\(M,S\)};\draw[->](3.45,1.53)--(4.3,1.53);\end{tikzpicture}\end{center}
Рис. $1$

\begin{center}\begin{tikzpicture}[task graph,>=stealth]\draw (0,3.5)--(0,0)--(2.25,0)--(2.25,3.5);\draw(-.1,0)--(2.4,0);\foreach \x in {0,.16,...,2.4}{\draw(\x,0)--++(-.15,-.16);}\fill(0,1.2)--(.23,1.2)--(0,.97)--cycle;\fill(2.25,1.2)--(2.02,1.2)--(2.25,.97)--cycle;\draw[fill=orange!30](0,2.35)rectangle++(2.25,.32);\draw(2.25,2.35)--(3.1,2.35);\draw[<->](2.85,0)--(2.85,2.35)node[midway,right]{\(H\)};\node at(1.2,3.15){\(p_0\)};\end{tikzpicture}\end{center}
Рис. $2$
}{
\smallskip\noindent\textbf{Решение}\par
Изобразим силы, действующие на поршень в процессе движения.
\begin{center}\begin{tikzpicture}[task graph,>=stealth]\draw[fill=orange!30](-.78,-.25)rectangle(.78,.25);\fill (0,0)circle(1.3pt);\draw[->](0,0)--(0,1.55)node[right]{\(p_2S\)};\draw[->](0,0)--(0,-.77)node[right]{\(Mg\)};\draw[->](0,0)--(0,-1.4)node[right]{\(p_0S\)};\end{tikzpicture}\end{center}
До момента, пока $p_2S < p_0S + Mg,$ поршень будет покоиться, давление газа увеличиваться. При достижении равновесия $p_2S=p_0S + Mg$ поршень начнет двигаться вверх, так как нагревание медленное, из-за чего процесс можно считать изобарным.
По первому закону термодинамики $Q=A + \Delta U.$
Изменение внутренней энергии равно $\Delta U= \frac{3}{2} \nu R\Delta T.$
Работа газа при изобарном расширении:

\[A=p_2\Delta V= ( p_0 + \frac{Mg}{S} ) S ( H - h ) = ( p_0S + Mg ) ( H - h ) .\]

Пусть начальный объем $V_0=Sh,$ конечный объем $V=SH.$
Запишем уравнение Клапейрона-Менделеева для каждого состояния газа:

$p_0V_0=\nu RT_0;$
$p_2V_=\nu RT_2.$

При объединении уравнений, получаем:

\[( p_0S + Mg ) H - p_0Sh=\nu R\Delta T.\]

Тогда изменение внутренней энергии газа равно:

\[\Delta U= \frac{3}{2} ( ( p_0S + Mg ) H - p_0Sh ) .\]

Тогда полученное газом количество теплоты равно:

$Q= \frac{3}{2} ( ( p_0S + Mg ) H - p_0Sh ) + ( p_0S + Mg ) ( H - h ) ;$
$Q= 2{,}5p_0S ( H - h ) + Mg ( 2{,}5H - h ) .$

Ответ: $Q= 2{,}5p_0S ( H - h ) + Mg ( 2{,}5H - h ) .$
}


% Source 7875: https://phys-ege.sdamgia.ru/problem?id=7875
\BankTask{7875}{
\begin{center}\begin{tikzpicture}[task graph,>=stealth]\draw(5,0)--(0,0)--(0,2.65)--(5,2.65);\draw[fill=orange!25](2.26,0)rectangle(2.66,2.65);\foreach \y in {.13,.3,...,2.63}{\draw(2.26,\y)--++(.4,-.12);}\draw(0,1.26)--(.18,1.26);\foreach \x in {.2,.45,...,1.95}{\draw (\x,1.26)arc(180:-145:.13 and .19);}\draw(2.05,1.26)--(2.26,1.26);\node at(1.05,1.91){\(k\)};\node at(.96,.55){\(p_0\)};\node at(4,1.15){\(p_0\)};\draw(2.63,2.1)--(3.8,3.14)node[right]{\(S\)};\draw(0,0)--(0,-.95)(2.26,0)--(2.26,-.95);\draw( .22,-.74)--(2.04,-.74)node[midway,above]{\(L\)};\end{tikzpicture}\end{center}
В горизонтальном цилиндре с гладкими стенками под массивным поршнем с площадью S находится одноатомный идеальный газ. Поршень соединен с основанием цилиндра пружиной с жесткостью k. В начальном состоянии расстояние между поршнем и основанием цилиндра равно L, а давление газа в цилиндре равно внешнему атмосферному давлению $p_{0}$ (см. рис.). Какое количество теплоты Q передано затем газу, если в результате поршень медленно переместился вправо на расстояние b?
}{
\smallskip\noindent\textbf{Решение}\par
Тепло, переданное газу, идет на изменение его внутренней энергии и на совершением им работы:

\[Q=\Delta U + A.\]

В начальном состоянии давление и объем газа равны $p_0$ и $V_0=SL.$ Следовательно, пружина была не деформирована. 
В конечном состоянии пружина была деформирована на величину b, возникла сила упругости $F_{\,\text{упр}}=kb,$ направленная влево. Запишем условие равновесия поршня в конечном состоянии:

\[p_0S + F_{\,\text{упр}}=p_{\,\text{к}}S.\]

Тогда давление воздуха в конечном состоянии $p_{\,\text{к}}=p_0 + \frac{kb}{S}$ и объем воздуха в конечном состоянии  $V_{\,\text{к}}=S ( L + b ) .$ 
Изменение внутренней энергии воздуха $\Delta U = \frac{3}{2} \nu R\Delta T.$ Тогда получаем:

\[\Delta U = \frac{3}{2} ( p_{\,\text{к}}V_{\,\text{к}} - p_0V_0 ) = \frac{3}{2} ( p_0Sb + kbL + kb ^2 ) .\]

Чтобы рассчитать работу, заметим, что в каждый момент времени, когда поршень сдвинут на x от начального положения давление равно $p=p_0 + \frac{kx}{S} =p_0 + \frac{k ( V - V_0 )}{S ^2} ,$ то есть давление линейно зависит от объема. Значит, на pV-диаграмме процесс расширения будет изображен отрезком прямой, а фигура под графиком будет являться трапецией, площадь которой равна:

\[A = \frac{p_0 + p_{\,\text{к}}}{2} ( V_{\,\text{к}} - V_0 ) =p_0Sb + \frac{kb ^2}{2} .\]

Таким образом, полученное воздухом количество теплоты равно:

\[Q = \frac{5}{2} p_0Sb + \frac{3}{2} kbL + 2kb ^2 .\]

Ответ: $Q = \frac{5}{2} p_0Sb + \frac{3}{2} kbL + 2kb ^2 .$
}

\end{document}
